\section{鲁、郑、宋：中原小国怎样活在大国之间}

洛邑以东、以南的道路把鲁、郑、宋连入同一片诸侯世界，却没有给它们同一种命运。鲁在曲阜，国力未必压过邻人，留下的年次却成为后人辨认春秋的经线；郑贴着王畿，既替周室办事，也最先感到王命不能再直接调动地方；宋夹在晋、楚之间，兵车往来时常是受压的一方，僵局形成时又可能成为双方都肯入席的地方。小国的作为不在于脱离大国，而在于看准道路、礼仪和资源暂时留下的缝隙。

这条经线不能向前无限延伸。约前768年，后出的鲁侯世系把\eventanchor{SA-0768-LU-HUI}{春秋·鲁惠公即位}安在东迁前后的继承序列中；它可提示鲁国的宗族政治已有前史，却不是曲阜保存下来的一册逐年档案。郑的情形更能显出空间先于史书的重要：传统叙事称郑武公先\eventanchor{SA-0760-ZHENG-HU}{春秋·郑武公伐胡}，继而\eventanchor{SA-0758-ZHENG-KUAI}{春秋·郑取郐}，取得面向中原道路的较稳据点。新郑一带城址能够帮助追问这种据点的空间条件，不能单独证成攻取的年份、协议或旧族如何安置。郑后来既近王畿又有自己的城邑腹地，正是在这种并不完整的前史中形成。

同时代的鲁历尚未开篇，其他小国也未因此没有继承问题。约前758年\eventanchor{SA-0757-WEI-WU-DEATH}{春秋·卫武公卒}，卫庄公承接黄河北岸的君位；约前729年\eventanchor{SA-0729-SONG-MU-ASCEND}{春秋·宋穆公即位}，宋在公族之间完成一次传世谱系所记的交接。两事的年数主要经《史记》及较晚传统整理，不能把它们写成可逐日复原的宫廷场景。不过它们提醒我们：鲁、郑、宋及其邻国进入《春秋》时，并非刚刚出现的国家，而是已有宗庙、支系和交通关系、却缺少同等密度编年的政治体。

\begin{event}{公元前722年}{鲁隐公元年与《春秋》年表开始}{可靠纪年}{鲁}\eventanchor{SA-0722-CHUNQIU}{春秋·鲁隐公元年}
曲阜的史官以“元年春王正月”落下第一笔。往后的字句短得近乎冷峻：会盟、来朝、伐国、葬君、日食，都被安进鲁君在位的年次。它用的是周王正月，站的却是鲁国朝廷的门槛；列国由此有了一条可以相互校验的时间线，后世也因此容易把鲁人看见的事误作整个天下唯一的视野。

这种可见性本身是一种政治资源。东迁后的王室尚能提供月令、爵名和礼仪的共同语言，却已不能替所有诸侯裁断争端。鲁国没有因保存年表便取得霸权；它所保住的，是把各国行动写成先后相承之事的权力。经文中看似平静的称谓、动词和取舍，正是后来的传家、经师和史家不断进入的门。

\term{史料边界}：隐公元年的经文是可靠的年代起点。《左传》《公羊传》《谷梁传》分别把这部短书展开为叙事、义例和经学解释；它们不能彼此自动证实。孔子删修《春秋》是影响深远的传统说法，但现存经文的档案基础、编纂层次和传抄过程仍是研究问题，不能把每一字都当作一位作者留下的现场笔录。
\end{event}

周正月给出了共同的年月和爵名，却不替各国消除边缘与城邑的麻烦。隐公二年，\eventanchor{SA-0721-LU-RONG-MEET}{春秋·鲁会戎于潜}把戎人写进鲁历最早一批外交记录；经文只证实会见，不足以推定戎人的组织、盟辞或归属。前718年，宋军\eventanchor{SA-0718-SONG-CHANGGE}{春秋·宋围长葛}，郑宋争执已落到围邑、守粮和道路上。前712年，郑伯\eventanchor{SA-0712-XU-FALL}{春秋·郑伯克许}，许君出奔，传世叙事又称郑将许地分给盟友和宗支；经年可定，分地比例和驻军安排却没有完整清册。礼仪的共同语言并未取消争夺土地、通道和可征发人力的压力。

鲁的历法也不能替鲁君解决继承。前712年，\eventanchor{SA-0712-LU-YIN-ASSASS}{春秋·鲁隐公遇弑}使代行国政的安排转为君位危机；弑君动机与羽父角色主要由后出的叙事塑形，最小可证的是隐公死而桓公继位。次年，\eventanchor{SA-0711-LU-HUAN-ASCEND}{春秋·鲁桓公即位}后即与郑会于垂、盟于越，并以许田牵连政治承认和土地权利。周的名分仍使会盟可以被说成合礼的关系，实际的承认却必须由相邻国家与可见资源来担保。

鲁史开篇的同一年，郑国已经显出另一种小国难题。郑庄公即位后，母族、弟弟共叔段和封邑之间的冲突终于公开。经文以“\eventanchor{SA-0722-ZHENG-DUAN}{春秋·郑伯克段于鄢}”记其事，留下的是宗君击败弟弟的结果；《左传》才铺开武姜偏爱、请封、筑城、庄公等待其变，以及“克段于鄢”的道德张力。可以确知的是，继承竞争和封邑所附的人力、城邑使郑的近畿权力不能自然集中；至于母子之间每一次言语、庄公是否预作圈套，属于后出叙事最浓的部分。

庄公取胜并没有使郑成为脱离周室的独立堡垒，反而使他更能以卿士身份同王室周旋。郑武公、庄公父子长期居于周朝政务的要位，郑既可调用本国的邑兵与粮食，也要为王室的衰弱承担距离最近的压力。平王晚年改倚虢公，周、郑互换人质，郑人又取温地之麦、成周之禾；《左传》把这些事连成一条逼向战争的线索。它提示冲突不只是抽象的“尊王”与“叛王”，还牵涉可耕之地、供军之粮和卿士职位；不过，具体交涉的先后和言辞仍须保留传文的叙事层次。

\begin{event}{公元前707年}{郑庄公与王师战于繻葛}{可靠纪年}{繻葛}\eventref{SA-0707-XUGE}{春秋·繻葛之战}
桓王终于率诸侯之师伐郑。繻葛的地望至今有不同考订，战场却足够靠近王畿，让这次交锋格外刺眼：王室以旧有的征发名义出兵，郑则以一国军政资源迎战。《春秋》只写“王师败绩于繻葛”，已把胜负钉在年表上；《左传》所写祝聃射中王肩、庄公战后处置，则把兵事编进礼与名分的戏剧性叙述，不能逐字当作战地实录。

郑军的胜利没有废去周王的名分。它显示的更为具体：当王室的兵源、财赋和指挥无法压倒近畿诸侯时，掌握城邑和军队的卿士可以拒命而仍保有位置。直接后果是王室军事威信受损；中期影响则是诸侯越来越依靠会盟、暂时从属和武力处理争端。郑的强势也有边界，它来自紧邻洛邑的特殊条件，不能无限外推为一种稳定霸权。
\end{event}

\begin{sourcenote}[郑庄公的家事与王事]
《春秋·隐公元年》与《春秋·桓公五年》分别固定了克段、王师败绩的年代骨架；《左传·隐公元年、三年及桓公五年》把封邑、质子、麦禾和箭伤编成连续故事，《史记·郑世家》又以汉人的叙述方式承接了这条线。互读可见郑的内政集中和周郑摩擦彼此相关，却不能把传文中完整的心理、对白和战术细节当作同时代材料已经独立证明的事实。
\end{sourcenote}

繻葛之后的郑也没有把近畿优势变成一条直线上升的国势。前701年庄公死，\eventanchor{SA-0701-ZHENG-ZHUANG-DEATH}{春秋·郑庄公卒}，昭公继位；次年祭仲逐昭公而立子突为厉公，\eventanchor{SA-0700-ZHENG-ZHAO-EXILE}{春秋·郑昭公出奔}把卿族立君公开推到国政中心。前697年厉公出奔、昭公复位，\eventanchor{SA-0697-ZHENG-ZHAO-RETURN}{春秋·郑昭公复位}；前695年昭公遇弑、子亹继立，\eventanchor{SA-0695-ZHENG-ZHAO-DEATH}{春秋·郑昭公遇弑}。这些年次能够彼此校核，祭仲受谁胁迫、谁在宫内谋画则主要属于《左传》《史记》保存的叙事。可以把握的结构性后果是：一国若以卿族、城邑和外部关系来处理继承，国君的外交承诺便可能随人事反复重置。前672年厉公死、文公继位，\eventanchor{SA-0672-ZHENG-LI-DEATH}{春秋·郑厉公卒}才给郑带来较长的君统延续；这不是内争已经消失，而是郑得以在齐、晋、楚之间继续回旋的必要条件。

鲁也不是只在竹简上观看别国。前694年桓公赴齐而死，\eventanchor{SA-0694-LU-HUAN-ASSASS}{春秋·鲁桓公死于齐}；齐鲁婚姻与会盟怎样转成危机，不能从后出叙事的私密情节倒推。前693年庄公即位，\eventanchor{SA-0693-LU-ZHUANG-ASCEND}{春秋·鲁庄公即位}，幼主朝廷须同时面对齐的压力和国内卿族的调度。前684年，齐军南来，鲁军在长勺迎战。\eventanchor{SA-0684-CHANGSHAO}{春秋·长勺之战}的经文确认齐伐鲁和鲁师获胜，说明曲阜仍能在边境保住自身；《左传》以曹刿论战的节奏解释“齐师败绩”，遂使鼓进鼓止成为后世最熟悉的春秋战场画面。那些言辞保存的是后人理解车战、军心与主将判断的方式，不等于我们能够据此复原每一步战术。鲁赢得的是一段边境喘息，而不是改写齐国动员能力的资本；齐的资源和联盟随后仍将转向更广阔的北方与中原。

宋的困局比鲁更裸露，也更能照见周礼名分之外的保护何其有限。前680年，齐在北杏会后伐宋，\eventanchor{SA-0680-QI-SONG}{春秋·齐伐宋}表明会盟背后仍须有惩罚力量；经文并不支持把它夸大为齐已取得稳定的宗主权。前660年狄人入卫，卫懿公死、旧都受重创，\eventanchor{SA-0660-WEI-FALL}{春秋·狄入卫}；传世叙事又称宋桓公接应卫遗民渡河、立戴公于曹地，并有齐军戍守，\eventanchor{SA-0660-WEI-DAI}{春秋·宋接卫遗民}。人数、馈赠和行军路线不能由此量化，但卫亡而复存说明小国的名号、人口和宗庙可以凭邻国接应暂时续接。前656年齐与江、黄伐陈，\eventanchor{SA-0656-CHEN-PUNITIVE}{春秋·齐伐陈}又显示夹在道路上的国家即使短期顺从，也要反复承受盟网的检验。

齐桓公死后，宋襄公试图把这样的礼仪记忆转成会盟声望。约前650年，\eventanchor{SA-0650-SONG-XIANG-ASCEND}{春秋·宋襄公争取诸侯领导}的具体即位年与会盟次序尚有异说；前639年盂地会盟中扣留诸侯，\eventanchor{SA-0639-SONG-BLINDNESS}{春秋·宋襄公扣留诸侯}则使协调关系带上强制色彩。它暴露出宋的难处：宋可以召集和陈说，却没有足以长期替代齐、晋的军政资源。前638年，他在泓水与楚军相遇。\eventanchor{SA-0638-HONG}{春秋·泓水之战}以宋师败绩收束，经文所能牢固确认的是战事、交战者和结果。后出的《左传》写宋襄公不肯趁楚军渡河、未成列时攻击，把一次败战塑成守礼而失机的寓言；《国语·宋语》与《史记·宋微子世家》又各自延续宋襄公的形象。与其把这个故事简单读作“仁义误国”，不如先看见宋所面对的资源差距：它想以礼仪号召诸侯，楚却能以北进的军政力量争夺同一批关系。泓水之后，宋的霸权尝试受挫，楚在中原的存在更难忽略，并接入了\eventref{SA-0632-CHENGPU}{春秋·城濮之战}前后的晋楚竞争。

泓水没有令宋退出战场。前634年楚围宋，\eventanchor{SA-0634-CHU-SONG-SIEGE}{春秋·楚围宋}，宋向晋、秦求援；围城的补给状况与使者行程主要见叙事，较稳的是宋的危机成为晋楚正面冲突前的直接触发。郑的保存也有相同的两面。前630年晋、秦围郑，秦后退兵，郑得以留下来，详见\eventref{SA-0630-QINJIN-ZHENG}{春秋·秦晋围郑}；前627年郑文公卒、穆公即位，\eventanchor{SA-0627-ZHENG-WEN-DEATH}{春秋·郑文公卒}，新君继承的不是安全，而是必须继续衡量秦、晋和楚的通道位置。前626年晋伐卫并取戚，\eventanchor{SA-0626-JIN-WEI-CAMPAIGN}{春秋·晋伐卫取戚}，朝聘义务已经可以落到边境土地和军队上。它们都说明“从属”并非一次选边，而是不断受兵、求援、交人和维持城邑的过程。

这种过程很快把宋、郑的内政也拉进来。约前611年宋昭公遇弑、文公继位，\eventanchor{SA-0611-SONG-ZHAO-ASSASS}{春秋·宋昭公遇弑}；弑君责任及晋国讨宋的说法主要来自较晚叙事，不能据此补出完整审判。前608年楚与郑侵陈、宋，晋赵盾救陈宋而又伐郑，\eventanchor{SA-0608-CHU-ZHENG-CHEN-SONG}{春秋·楚郑侵陈宋与晋救援}，一条陈宋粮道和郑的地理位置足以把数国卷入，却不意味着已有统一战区指挥。次年宋华元在大棘败而被俘，\eventanchor{SA-0607-SONG-ZHENG-DAJI}{春秋·宋郑战于大棘}；《左传》的军中饮食故事不能替代战术纪录，战败、俘虏和随后城防压力则是宋不得不承受的现实。

郑亦无法把外交同宫廷切开。前604年郑灵公被子公所杀、襄公继立，\eventanchor{SA-0604-ZHENG-LING-ASSASSINATION}{春秋·郑灵公遇弑}；食鼋等细节属于传文的叙事层。前601年楚庄王围郑，\eventanchor{SA-0601-CHU-ZHENG}{春秋·楚庄王围郑}，郑在晋楚间调整从属，恰说明一座近畿城邦的摇摆足以成为两强争夺中原的杠杆。前595年楚再围宋，\eventanchor{SA-0595-CHU-SONG-SIEGE}{春秋·楚围宋与宋守城}；宋以城守、使者和赎和维持国祀，围困时长和粮食数字却不应由传文作出精确推算。小国没有因为守住都城便免除代价，防务和外交的成本只会更长久地进入城市社会。

鲁的存续同样伴随资源外流。前609至608年，文公卒，宣公在东门氏支持下即位，\eventanchor{SA-0609-LU-XUAN-ASCEND}{春秋·鲁宣公即位}；东门氏在继承中的实际作用须和《左传》的叙事分开。齐随后取鲁济西田，\eventanchor{SA-0608-QI-JIXI}{春秋·齐取鲁济西田}，鲁以土地和礼聘换取对新君的认可，田界细目和单一书面交易均不可确指。前590年成公即位，季孙行父等继续主政，\eventanchor{SA-0590-LU-CHENG-ASCEND}{春秋·鲁成公即位与季孙行父执政}；这位季孙并非抽象的“大夫势力”，而是把职位、邑兵和对外事务连在同一户族网络中的执政者。鲁还可以随盟主出兵，却越来越难由公室单独决定军政资源。

邻国的经验补足了这层脆弱性。前589年齐伐卫，卫在新筑败而向晋求援，\eventanchor{SA-0589-WEI-XINZHU}{春秋·卫败于新筑}；参战兵力不能从传文的勇怯评语中算出。前582年宋共公卒、平公继位，\eventanchor{SA-0582-SONG-GONG-DEATH}{春秋·宋平公即位}，宋仍在两强间保有都城与会盟地点，却也必须把这种位置转化为可以降低风险的交涉。

宋人并未因此只能挨打。晋、楚连年争胜，宋恰在军队和使者必须经过的中间地带；降低战争风险首先是它自己的生存需要。前579年，宋大夫华元推动晋楚在宋结盟。\eventanchor{SA-0579-HUAYUAN}{春秋·宋华元调停}所见的，是一次并未终结竞争的停歇：两强都还保有兵力和附属国，也都暂时不愿继续承受征发的代价。《春秋·成公十二年》提供会盟的年次，《左传》和《国语·宋语》所存谈判言辞则带有很强的外交修辞。华元未能创造一个凌驾诸侯的仲裁者，却让“由宋居间、以会盟减损战事”成为可再使用的办法。

不过，华元的会盟并没有拆除争夺的网络。前572年，晋联合宋、卫、曹、莒、邾、滕、薛围彭城以处置宋国叛臣，\eventanchor{SA-0572-PENGCHENG}{春秋·晋诸侯围彭城}；“为宋讨”是盟军的立场，彭城居民与归属并不能由此一语定论。同年晋军入郑郛、败郑徒兵于洧上，\eventanchor{SA-0572-JIN-ZHENG-WEISHANG}{春秋·晋军入郑郛}，楚子辛救郑又侵宋吕、留，郑亦取宋犬丘，\eventanchor{SA-0572-CHU-SONG-RESCUE}{春秋·楚郑侵宋以救郑}。三件事并列在年表中，最能显示宋、郑不是会盟图上的静点，而是迫使两大集团分散兵力的城邑与边地。

郑的继承于是越来越依赖多卿分掌。前571年成公卒、僖公即位，子罕、子驷、子国分掌国政和军职，\eventanchor{SA-0571-ZHENG-CHENG-DEATH}{春秋·郑僖公即位与诸卿分掌}；这不是一张可以精确还原权限的官制图，却可见新君须借既有执政网络应对外压。前570年晋伐许，\eventanchor{SA-0570-JIN-XU-CAMPAIGN}{春秋·晋伐许}，又把郑楚关系的压力推向较小的许国。前566年僖公卒、简公继位，\eventanchor{SA-0566-ZHENG-JIAN-ASCEND}{春秋·郑简公即位}；前564年晋以郑“不祀”晋文公为由伐郑，\eventanchor{SA-0564-JIN-ZHENG-CAMPAIGN}{春秋·晋以不祀伐郑}。祭祀理由究竟是实在义务还是出兵语汇，仍须讨论；可确定的是，礼仪能够为惩罚命名，却不能免去郑的城邑与民力。前562年晋悼公又率诸侯伐郑，\eventanchor{SA-0562-JIN-WEI-ATTACK}{春秋·晋悼公伐郑}，反复兵役正是各方终于愿意考虑弭兵的现实背景。

\begin{event}{公元前546年}{宋向戌调停晋楚于弭兵之会}{可靠纪年}{宋}\eventref{SA-0546-MIBING}{春秋·弭兵之会}
经过\eventref{SA-0575-YANLING}{春秋·鄢陵之战}，晋、楚都还能够动员，谁也难以用一场战役迫使对方退出中原。宋的执政大夫向戌在这道僵局中来往于两边；《左传·襄公二十七年》用“将平晋、楚”概括他的目标。调停者不是站在冲突之外的裁判，而是必须让两方相信：把会面放在宋，所得会比继续出兵更多。

会盟在宋举行，所达成的并不是现代意义的和平条约。晋、楚各自的从属网络被暂时并列承认，诸侯因而要承受双重而紧张的关系；齐、秦等国也没有因此纳入一个统一的指挥体系。宋获得了罕有的外交能见度，同时要承担接待、传递意向和维持中立的风险。它的力量不在替晋楚下令，而在于把双方不愿放弃的礼仪空间变成一张可以谈判的席位。

弭兵减少了晋楚直接决战，不能消除兼并，更没有设立一支能执行盟约的常备力量。外部冲突暂缓后，诸国卿族和执政集团反而有更多时间经营封邑、家臣和军队。它是消耗之后的均势安排，不是春秋秩序的终局；华元的先例和向戌的成功，恰好说明宋的主动何其依赖晋楚谁也不能独胜的时刻。
\end{event}

\begin{causalchain}[从王师败绩到宋地会盟]
郑在\eventref{SA-0707-XUGE}{春秋·繻葛之战}所揭示的，是王室已难独断近畿；到\eventref{SA-0546-MIBING}{春秋·弭兵之会}，问题已不再是恢复王命，而是怎样让并立的晋楚减少相互耗损。这不是从一场败仗直线推出一场会盟，而是旧有裁断机制逐渐失效后，各地分别以武力、盟约和调停填补空缺的过程。
\end{causalchain}

郑在这种两强夹缝中走得更久，代价是把许多政治争执带进城内。前543年前后，子产开始主掌郑国政务。\eventanchor{SA-0543-ZICHAN-POWER}{春秋·子产执政}并非一纸新制的突然降临：郑要在晋楚之间选择进退，又要处理卿族、城邑和诉讼，稳定可预期的行政比一时扩张更迫切。《左传·襄公三十年》《国语·郑语》与《史记·郑世家》所保存的子产言行，多已成为后世的政治记忆；可以用来观察小国治理的问题，不能把其中每段议论视为当日的速记。

前541年，传世材料又把子产与田制、赋役、城防的整顿连在一起，\eventanchor{SA-0541-ZICHAN-LAND}{春秋·子产整顿田制赋役与城防}。事迹的方向可见，具体法令的年次、范围和受益者却有争议；它不应被改写成一套先于战国的完备财政制度。较稳妥的理解是，郑若要在大国兵役与城防压力之间活下去，就须让征调和裁断少一些临时性的任意。这也解释了为什么子产的公共秩序实践既关乎国内卿族，也关乎小国对外还能否持续回应。

前536年，郑人铸刑书，\eventanchor{SA-0536-ZHENG-XINGSHU}{春秋·子产铸刑书}把原先更多依赖贵族议礼的裁断规则以器物公布。《左传·昭公六年》还保存叔向致书批评的说法，使后世常把它读成礼与法的一次正面冲突。实际公布的条文、适用范围，以及书信的真伪和成文层次，学界仍有讨论；但这一举动至少表明，法律知识不再完全锁在少数人对礼例的掌握之中。它提高了裁断的可预期性，也可能使原有贵族对解释规则的垄断受到挑战。

郑的秩序努力并不只面对郑人。前533年许迁于夷，\eventanchor{SA-0533-XU-YI}{春秋·许迁于夷}；楚主导迁徙的程度、夷地位置和迁民数量都有历史地理上的争论，但许的宗庙要以更深嵌入楚方势力为代价才得延续。子产去世前后的郑，正是在这种不断移动、受兵又求存的邻国世界中处理自己的田亩、城墙和案件。

前522年子产去世，\eventanchor{SA-0522-ZICHAN-DEATH}{春秋·子产去世}后的郑仍须在晋楚之间维持位置。传世材料中民众哀悼、后继者称誉子产的篇幅很动人，却是《左传》《国语》以政治理想回望一位执政者的方式。较能把握的是他的长期经营没有替郑消去外压，只是使这座近畿国家在资源有限时拥有了较稳的处置手段；后人所称道的子产，也由此成为评价小国政治的尺度，而不是一段可以脱离条件复制的答案。

鲁的“记录者”位置同样没有保护鲁君。前542年襄公卒、昭公即位，\eventanchor{SA-0542-LU-ZHAO-ASCEND}{春秋·鲁昭公即位}时，季、孟、叔孙三家已掌握邑兵和执政职位；后来的冲突不能倒推为当年每一家已有同一套固定权力，却足以说明新君没有独立军政基础。前517年，\eventanchor{SA-0517-LU-ZHAO-EXILE}{春秋·鲁昭公出奔}的经文记下君主奔齐，鲁国公室的空心化遂无法再用礼仪遮蔽。《左传》关于攻郓等战事的展开，能让我们看见公室与季氏集团如何在城邑间争夺，细部仍属叙事材料，需要和经文的结果分开读。国君出奔并非鲁政一夕崩解，而是多年资源由公室流向卿族的公开显影。

昭公没有回到曲阜。前509年，他卒于乾侯，定公即位，\eventanchor{SA-0509-LU-DING-ASCEND}{春秋·鲁定公即位}结束了名义上的空位，却没有收回三桓掌握的邑兵和外交权。前504年郑灭许、俘许男，\eventanchor{SA-0504-ZHENG-XU}{春秋·郑灭许}；许此前多次迁徙而城防衰弱，郑的出兵实际控制者和许都位置仍须辨析，但这次失国说明春秋末的小国更难只凭会盟名分保存自己。鲁的内部危机与郑的邻近扩张并行，不能彼此互作原因，却共同显示资源网络已比旧封国的名号更能决定行动边界。

十五年后，季氏的家臣阳虎也在斗争中失势出奔。经文记下\eventanchor{SA-0502-YANGHU}{春秋·阳虎出奔}，让权力再向下移的一刻浮出纸面：公室受制于三桓，三桓又可能受制于掌兵掌邑的家臣。阳虎的动机及其同孔子的后出关联，不能以传说补成确定事实；但他的败走并未把权力送回鲁君，反而显示封邑、家臣和武装一旦形成网络，谁控制它们就可以暂时越过名分行事。

鲁廷仍会尝试把礼仪重新用作外交和内政的杠杆。前500年齐鲁会于夹谷，孔子以相礼身份随定公赴盟，\eventanchor{SA-0500-JIAGU}{春秋·齐鲁夹谷之会}；临场言辞与齐归田之间的因果链主要由《左传》《史记》组织，经文所可凭的是会盟本身。齐归还郓、讙、龟阴等地的记载，至少显示鲁公室一度能把礼仪、边地和外部均势放在同一次交涉中。前498年叔孙、季孙先后堕郈、费，鲁君围成而未克，\eventanchor{SA-0498-DUO-SANDU}{春秋·鲁堕三都}；后出的传记放大了孔子的决策作用，较稳的事实是三桓据兵征赋的都邑并未被一举解除，公室的动员边界反而显露出来。

同一时期，较小的蔡为了延续宗庙迁到州来，\eventanchor{SA-0493-CAI-ZHOULAI}{春秋·蔡迁州来}。楚围蔡后的迁民数、旧城利用和谁主导迁徙都难详考，但蔡对吴的保护更深依赖，重画了吴楚之间的淮河通道。鲁在夹谷的一次收获和蔡的迁徙不是一条单线的胜败叙事：前者是利用礼仪与外部关系争取余地，后者则以地理迁移换取存续；两种做法都没有恢复一个能独断诸侯的周王室。

\begin{event}{公元前481年}{获麟与齐国陈氏的权力跃升}{可靠纪年}{鲁与齐}\eventanchor{SA-0481-LIN}{春秋·获麟与陈氏政变}
鲁哀公十四年，鲁国经文写下“西狩获麟”。四字像一扇忽然合上的门，后世遂常以\eventref{SA-0481-END}{春秋·记事终点}标识《春秋》的终篇。它首先是文本传承的终点，不是列国政治在这一年忽然换了一副面貌。

同年，齐国陈恒杀齐简公壬于舒州，改立其弟骜。经文称“陈恒”，后世通称“田恒”，是同一支陈、田氏族的不同书写；这次行动使田氏在齐国内部的控制越过又一道门槛。它仍不是\eventref{WQ-0386-TIANQI}{战国·田氏代齐}的提前完成。若把前481年直接叫作“田氏代齐”，便会抹去贵族家族长期积累封邑、家臣和军政资源的过程，也会把政变后的实际控制误作法统上的改朝换代。

获麟与陈恒弑君同年出现，并不意味着前者造成后者。鲁的记录停止，齐的卿族上升，却让人同时看见春秋末叶两种转折：一边是旧编年传统给出的收束，一边是国君与世卿之间的权力天平继续倾斜。此前弭兵保留的均势并未冻结各国内部的竞争；齐国的变化，正是这种长期重组在一国之内最显眼的一幕。
\end{event}

\begin{textwindow}[《春秋》的最后一笔]
“西狩获麟。”经文只有四字。《公羊传》把它解释为孔子“吾道穷矣”的象征；《左传》与《史记》的齐事则把注意力转向陈恒弑君及其后的权力安排。原文无需逐字翻译；需要分开的是经文的停止、经学传统赋予它的象征，以及齐国内部实际而漫长的家族权力转移。
\end{textwindow}

\begin{turningpoint}[制度得失：小国的可见性与脆弱性]
鲁的年表、郑的近畿军政与宋的调停席位，都是小国在制度缝隙中取得的能力。它们分别解决记录、拒压和沟通的当下难题，成本却是高度依赖王室名分、晋楚均势和贵族网络。季孙行父主政、三桓与鲁君相持、阳虎以家臣身份一度越出主家，都不等于一套整齐的“贵族政体”，却显示职位、邑兵和家内网络怎样逐层离开公室的直接控制；郑以子产的征调与刑书求取可预期性，也不能消去强邻和卿族的压力。到鲁君出奔、家臣用事，以及陈、田氏越过齐君而行动时，封邑、家臣和军政资源已在改变权力的归属：礼仪仍能命名秩序，却越来越难单独约束掌握资源的家族。这正是春秋晚期通向更深制度重组的条件，而非某一年骤然开始的战国。
\end{turningpoint}
